% Use this documentclass for the PNAS style:
\documentclass[9pt,twocolumn,twoside,lineno]{pnas-new}

% This is for NeurIPS style:
%\documentclass{article}


\makeatletter
\@ifclassloaded{pnas-new}
{
  % Use the lineno option to display guide line numbers if required.

\templatetype{pnasresearcharticle} % Choose template
% {pnasresearcharticle} = Template for a two-column research article
% {pnasmathematics} %= Template for a one-column mathematics article
% {pnasinvited} %= Template for a PNAS invited submission

\setcounter{secnumdepth}{0}

\input{pre-macros}

\title{\input{title.tex}}

% Use letters for affiliations, numbers to show equal authorship (if applicable) and to indicate the corresponding author
\author[a,1]{Atticus Geiger}
\author[b]{Alexandra Carstensen}
\author[b]{Michael C.~Frank}
\author[a]{Christopher Potts}

\affil[a]{Department of Linguistics, Stanford University}
\affil[b]{Department of Psychology, Stanford University}

% Please give the surname of the lead author for the running footer
\leadauthor{Geiger}

% Please add here a significance statement to explain the relevance of your work
\significancestatement{The equality relation is a key case study for
broader questions of abstract relational reasoning. Earlier work concluded that neural networks are incapable of learning general solutions
to equality-based reasoning tasks, but this claim is in tension with the extraordinary recent
success of neural networks in artificial intelligence. We resolve this tension
by demonstrating these previous negative results hold only when neural networks are
provided hand-crafted feature representations. When neural networks are instead provided the random or pretrained
 representations that are now the norm in state-of-the-art work,
 these models easily learn general solutions to equality-based reasoning tasks,
 suggesting that essential aspects of symbolic reasoning can emerge from non-symbolic learning processes.}

% Please include corresponding author, author contribution and author declaration information
\authorcontributions{All four authors contributed to paper writing and experimental design. Atticus Geiger and Chris Potts wrote the code that was used to run the experiments.}
\authordeclaration{The authors declare no conflicts of interest.}
% \equalauthors{\textsuperscript{1}A.O.(Author One) and A.T. (Author Two) contributed equally to this work (remove if not applicable).}
\correspondingauthor{\textsuperscript{1}To whom correspondence should be addressed. E-mail: atticusg@stanford.edu}

% Keywords are not mandatory, but authors are strongly encouraged to provide them. If provided, please include two to five keywords, separated by the pipe symbol, e.g:
\keywords{relational reasoning $|$ neural networks $|$ generalization}

\begin{abstract}
  Humans have a remarkable capacity to reason about abstract relational structures, an ability that may support some of the most impressive, human-unique cognitive feats. Because equality (identity) is a simple and ubiquitous relational operator, equality reasoning has been a key case study for the broader question of abstract relational reasoning. This paper revisits the question of whether equality can be learned by neural networks that do not encode explicit symbolic structure. Earlier work arrived at a negative answer to this question, but that result holds only for a particular class of hand-crafted feature representations. In our experiments, we assess out-of-sample generalization of equality using both arbitrary representations and representations that have been pretrained on separate tasks to imbue them with  structure. In this setting, even simple neural networks are able to learn basic equality with relatively little training data. In a second case study, we show that sequential equality problems (learning ABA-patterned sequences) can be solved with only positive training instances. Finally, we consider a more complex, hierarchical equality problem, but this challenge requires vastly more data. However, using a pretrained equality network as a modular component of this larger task leads to good performance with no task-specific training ("zero-shot" generalization). Overall, these findings indicate that neural models are able to solve equality-based reasoning tasks, suggesting that essential aspects of symbolic reasoning can emerge from data-driven, non-symbolic learning processes.

\end{abstract}

\dates{This manuscript was compiled on \today}
\doi{\url{www.pnas.org/cgi/doi/10.1073/pnas.XXXXXXXXXX}}

\begin{document}

\maketitle
\thispagestyle{firststyle}
\ifthenelse{\boolean{shortarticle}}{\ifthenelse{\boolean{singlecolumn}}{\abscontentformatted}{\abscontent}}{}

% If your first paragraph (i.e. with the \dropcap) contains a list environment (quote, quotation, theorem, definition, enumerate, itemize...), the line after the list may have some extra indentation. If this is the case, add \parshape=0 to the end of the list environment.

% \section{Introduction}\label{sec:introduction}

  \newcommand{\papermode}{pnas}
}
{
  \input{pre-arxiv}
  \newcommand{\papermode}{arxiv}
}
\makeatother

\newcommand{\newatticus}[1]{{\color{red}#1}}
\newcommand{\newmike}[1]{{\color{blue}#1}}
\newcommand{\methref}{Methods}
\newcommand{\SIref}{SI}

\dropcap{O}ne of the key components of human intelligence is our ability to reason about abstract relations between stimuli. Many of the most unremarkable human activities -- scheduling a meeting, following traffic signs, assembling furniture -- require a fluency with abstraction and relational reasoning that is unmatched in nonhuman animals. An influential perspective on human uniqueness holds that relational concepts are critical to higher-order cognition \citep[e.g.,][]{Gentner:2003}. By far the most common case study of abstract relations has been equality.\footnote{We use the term ``equality'' here, though different literatures have also used ``identity.''} Equality is a valuable case study because it is simple and ubiquitous, but also completely abstract in the sense that it can be evaluated regardless of the identity of the stimuli being judged.

Equality reasoning has been studied extensively across a host of systems and tasks, with wildly variant conclusions. In some studies, equality is very challenging to learn: only great apes with either extensive language experience or specialized training succeed in matching tasks in which a \emph{same} pair, AA, must be matched to a novel same pair, BB \cite{Premack:1983,thompson:2001}. Preschool children also struggle to learn these regularities in a seemingly similar task \citep{walker:2016}. In contrast, other studies suggest that equality is simple: bees are able to learn abstract identity relationships from only a small set of training trials \cite{giurfa:2001,avargues:2011}, and human infants can generalize identity patterns \cite{anderson:2018} and succeed in relational matching tasks \cite{ferry:2015}. We take the central challenge of this literature to be characterizing the conditions that lead to success or failure in learning an abstract relation in a way that can be productively generalized to new stimuli \citep{carstensen:inpress}.

The learning task in all of these cases can be described using the predicate \emph{same} (or equivalently, =), which operates over two inputs and returns {\sc true} if they are identical in some respect. One perspective in the literature is that success in these learning tasks implies the presence of an equivalent symbolic description in the mind of the solver \cite{marcus:1999,Premack:1983}. This view does not provide a lever to distinguish which of these tasks are trivial and which are difficult, however. Further, it can fall prey to circularity: because newborns show sensitivity to identity relations \cite{gervain:2012}, then it would follow from this argument that they must have symbolic representations. If this logic applies also to bees, then we presuppose symbolic representations universally and have no account of the gradient difficulty of different tasks for different species.

An explanation of when same--different tasks are trivial and when they are difficult requires a theoretical framework beyond the symbolic/non-symbolic distinction. To make quantitative predictions about task performance, such a framework should ideally be instantiated in a computational model that takes in training data and learns a solution that generalizes when assessed with stimuli analogous to those used in experimental assessments. Symbolic computational models \cite[e.g.,][]{frank:2011} can be used to make contact with data about the breadth of generalization, but they require the existence of a symbolic equality predicate and hence again presuppose symbolic abilities in every case of success. Ideally, we would want a model that describes under what conditions \emph{same} is easy and under what conditions it is hard or unlearnable -- and how learning proceeds in hard cases. Here, we aim to lay the foundation for the development of such an account.

We are inspired by an emergent perspective in the animal learning literature that the representations underlying non-human animals' and human infants' successes in equality reasoning tasks are graded \cite{wasserman:2017}. This view acknowledges the increasing evidence that other species like pigeons \cite{cook2007}, crows \cite{smirnova2015}, and baboons \cite{fagot2011} can make true, out-of-sample generalizations of \emph{same} and \emph{different} relations, but it also recognizes that the observed patterns of behavior do not show the hallmarks of all-or-none symbolic representations. Instead, performance is graded. Out-of-sample generalization is possible but the level of performance depends critically on the diversity of the training stimuli \cite[e.g.,][]{castro2010}. Success requires hundreds, thousands, or even tens of thousands of training trials. And the outcome of learning is noisy and imperfect.

\newmike{These learning signatures appear to be a close match to the kind of learning exhibited by neural networks. Such networks are a flexible framework for arbitrary function learning, which have enjoyed a huge resurgence of interest in recent years in the fields of artificial intelligence, neuroscience, and cognitive science \cite[e.g.,][]{lecun2015,saxe2019}.

Marcus et al. \cite{marcus:1999} argued that a broad class of recurrent neural networks were unable to learn equality relations, however. These claims were  challenged by the presentation of evidence that some forms of neural networks are able to learn (at least aspects of) Marcus et al.'s equality tasks  \cite{dienes:1999,seidenberg:1999a,seidenberg:1999b,elman:1999,negishi:1999}, yet these examples were not uncontroversial. The subsequent debate \cite[reviewed in][]{alhama:2019} revealed a striking lack of consensus on some of the ground rules regarding what sort of generalization would be required to show that the learned function was suitably abstract.

Neural network models have recently been extraordinarily succesful on tasks such as natural language inference \cite{Bowman:2015, Williams:2018}, question answering \cite{Rajpurkar:2016, Rajpurkar:2018}, or visual reasoning \cite{Johnson:2017}, all of which are far more complex than equality-based tasks. In light of these findings, it may be surprising that the debate over equality-based reasoning is unresolved \cite{alhama:2019}. Yet, despite these developments, recent work on equality-based reasoning tasks takes as its starting point the conclusion that neural networks are unable to succeed using standard architectures and general purpose learning algorithms \cite{alhama2018,weyde2019,weyde2018,kopparti2020}. Similar negative results have been also found in the domain of computer vision \cite{Kim:2018,Fleuret:2011,Bengio:2016}, though the equality-based reasoning tasks in these experiments are typically more complex. Buttressing these claims, analyses of relational reasoning more broadly suggest that standard neural network architectures are not suitable for these tasks \cite{santoro2018}.

In sum, there is a tension between the success of modern deep learning models  and the claim that neural networks cannot learn equality-based reasoning tasks.
However, we suspect this tension exists because only a narrow range of network architectures and representations was explored in this earlier literature, in part because it predated many important innovations in neural network design. Thus we revisit the debate here, using a broader range of architectures and representations and adopting stringent criteria for generalization. In particular we explore random and pretrained representations, which have facilitated some of the extraordinary successes of modern artificial intelligence \cite{CollobertWeston:2011,Mikolov-etal:2013,pennington-socher-manning:2014:EMNLP2014,Peters-etal:2018,Devlin-etal:2019}. The use of pretrained representations to solve downstream tasks is argued to be a hallmark of natural learning systems \cite{zador2019}, has been an important feature of historical models from cognitive science \cite[e.g.,][]{landauer1997,macrae1993}, and is essential in the latest wave of state-of-the-art natural language processing models \cite{Devlin-etal:2019,Liu:2019,Radford:2019,Brown:2020}.
% While there has been criticisms of the datasets and methods used to evaluate neural models \cite{jia:2017,Kim:2018,nie:2020}

% We adopt stringent criteria for generalization and considering a broader range of representations across three models.
We model three cases of equality-based reasoning that have featured prominently in discussions of the role of symbols in relational reasoning (\figref{fig:tasks}): (1) learning to discriminate pairs of objects that exemplify the relation \emph{same} or \emph{different}, (2) learning sequences with repeated \emph{same} elements \cite{marcus:1999}, and (3) learning to distinguish hierarchical \emph{same} and \emph{different} relations in a context with pairs of pairs exemplifying these relations \cite{Premack:1983}. Across these three models, we find strong support for their ability to learn equality relations. These results should serve to revise the conclusions of the earlier debate.

Marcus and colleagues \cite{marcus:1999,marcus:2001} showed experimentally that neural networks using feature representations cannot generalize to binary features unseen in training. We agree with this claim (and support it with a direct mathematical argument in \appref{app:generalization}). However, they concluded from this result that neural networks will need to have primitive symbolic operators to solve hard relational reasoning tasks, \newatticus{which is a solution that has been pursued in recent machine learning research \cite{weyde2019,weyde2018, kopparti2020}}. On this point, we disagree. Our experiments show that networks without these primitives can solve a range of these tasks, as long as they use \newatticus{random or pretrained} representations. Overall, these findings suggest that some essential aspects of symbolic reasoning can emerge from entirely data-driven, non-symbolic learning processes.

Our work here makes three contributions. First, we resolve this longstanding debate by demonstrating neural networks are able to learn  equality relations when provided with pretrained or random representations. Second, we modify the standard architecture of a recurrent neural network to allow it to learn the sequential equality task with no negative feedback. Negative evidence was dismissed as an unreasonably strong learning regime in the original debate over these issues \citep[e.g.,][]{marcus:1999a}, and we show that this learning regime is not necessary. Third, we show that a model pretrained on the simple equality task can achieve zero shot generalization to the hierarchical equality task, suggesting that pretraining might provide an account of how some organisms succeed on hard relational learning tasks. In sum, we believe these three contributions represent significant progress in our understanding of neural networks' ability to perform equality-based reasoning.}

\begin{figure}[tp]
\centering
\includegraphics[scale=0.20]{modelfigv1.pdf}
\caption{Relational reasoning tasks. Green and red mark positive and negative training examples, respectively. The sequential task (Model~2) uses only positive instances, and a model succeeds if, prompted with $\alpha$, it produces a sequence $\beta \ \alpha$ where $\beta \neq \alpha$. For the hierarchical task (Model~3), we show that a model trained on the basic task (Model~1) is effective with no additional training.}
\label{fig:tasks}
\end{figure}

\newcommand{\mysquare}{
\begin{tikzpicture}[scale=0.5]
 \node[rectangle, fill=red!100, minimum height=3mm, minimum width=3mm]{};
\end{tikzpicture}}

\newcommand{\bluetriangle}{

\begin{tikzpicture}[scale=0.25]
 \node[regular polygon,regular polygon sides=3, fill=blue!100, inner sep=2pt]{};
\end{tikzpicture}}
\newcommand{\redtriangle}{
\begin{tikzpicture}[scale=0.5]
 \node[regular polygon,regular polygon sides=3, fill=red!100, inner sep=2pt]{};
\end{tikzpicture}}
\newcommand{\bluesquare}{
\begin{tikzpicture}[scale=0.5]
 \node[regular polygon,regular polygon sides=4, fill=blue!100, minimum height=3mm, minimum width=3mm]{};
\end{tikzpicture}}
\newcommand{\redsquare}{
\begin{tikzpicture}[scale=0.5]
 \node[regular polygon,regular polygon sides=4, fill=red!100, minimum height=3mm, minimum width=3mm]{};
\end{tikzpicture}}
\newcommand{\bluepentagon}{
\begin{tikzpicture}[scale=0.5]
 \node[regular polygon,regular polygon sides=5, fill=blue!100, minimum height=3mm, minimum width=3mm]{};
\end{tikzpicture}}
\newcommand{\redpentagon}{
\begin{tikzpicture}[scale=0.5]
 \node[regular polygon,regular polygon sides=5, fill=red!100, minimum height=3mm, minimum width=3mm]{};
\end{tikzpicture}}

\newcommand{\myrectangle}{
\begin{tikzpicture}[scale=0.25]
 \node[rectangle, fill=blue!100, minimum height=3mm, minimum width=5mm]{};
\end{tikzpicture}}

\newcommand{\mypentagon}{
\begin{tikzpicture}[scale=0.5]
 \node[regular polygon,regular polygon sides=5, minimum height=4mm, fill=red!100]{};
\end{tikzpicture}}


\newcommand{\mypolygon}{
\begin{tikzpicture}[scale=0.5]
 \node[regular polygon,regular polygon sides=7, minimum height=4mm, fill=blue!100]{};
\end{tikzpicture}}

\newcommand{\minshape}{6mm}

\begin{figure*}[t!]
\setlength{\arraycolsep}{2pt}
\centering
\small
\begin{subfigure}[t]{0.02\textwidth}
 \centering
 \begin{tabular}[b]{@{} c @{}}
  \midrule
  \bluetriangle \\
  \redtriangle \\
  \bluesquare \\
  \redsquare \\
  % \bluepentagon \\
   $\vdots$ \\
  \redpentagon \\
  \bottomrule
 \end{tabular}
 \caption{}
 \label{fig:entities}
\end{subfigure}
\begin{subfigure}[t]{0.16\textwidth}
 \centering
 \renewcommand{\arraystretch}{1.18}
 $\begin{array}[b]{*{6}{c}}
   \toprule
   \bluetriangle & \redtriangle & \bluesquare & \redsquare & \bluepentagon & \redpentagon \\
   \midrule
   1 & 0 & 0 & 0& 0& 0 \\
   0 & 1 & 0 & 0& 0& 0 \\
   0 & 0 & 1 & 0& 0& 0 \\
   0 & 0 & 0 & 1& 0& 0 \\
   % 0 & 0 & 0 & 0& 1&0 \\
   \multicolumn{6}{c}{\vdots} \\
   0 & 0 & 0 & 0& 0& 1 \\
  \bottomrule
 \end{array}$
 \caption{Localist.}
 \label{fig:reps:localist}
\end{subfigure}
\begin{subfigure}[t]{0.09\textwidth}
 \centering
 \renewcommand{\arraystretch}{1.18}
 $\begin{array}[b]{*{3}{r}}
   \toprule
   \rotatebox{90}{red} & \rotatebox{90}{blue} & \rotatebox{90}{sides}\\
   \midrule
   0 & 1 & 3  \\
   1 & 0 & 3  \\
   0 & 1 & 4  \\
   1 & 0 & 4 \\
   % 0 & 1 & 5 \\
   \multicolumn{3}{c}{\vdots} \\
   1 & 0 & 5 \\
   \bottomrule
 \end{array}$
 \caption{\newatticus{Property}.}
 \label{fig:reps:symbolic}
\end{subfigure}
\begin{subfigure}[t]{0.082\textwidth}
 \centering
 \renewcommand{\arraystretch}{1.18}
 $\begin{array}[b]{rr}
   \toprule
   d_{1} & d_{2} \\
   \midrule
   -0.4 &  0.2 \\
   0.2 &  -0.2 \\
   0.4 &  0.0 \\
   -0.2 &  -0.0 \\
   % 0.2 &  0.4 \\
   \multicolumn{2}{c}{\vdots} \\
   0.5 &  0.2 \\
   \bottomrule
 \end{array}$
 \caption{Random.}
 \label{fig:reps:random}
\end{subfigure}
\begin{subfigure}[t]{0.09\textwidth}
 \centering
 \renewcommand{\arraystretch}{1.18}
 $\begin{array}[b]{rr}
   \toprule
   d_{1} & d_{2} \\
   \midrule
   -1.3 &  0.4 \\
   -0.1 &  -1.1 \\
   -0.9 &  1.2 \\
   0.6 &  -0.6 \\
   %-0.9 &  2.1 \\
   \multicolumn{2}{c}{\vdots} \\
   1.2 &  -0.1 \\
   \bottomrule
 \end{array}$
 \caption{Pretrained.}
 \label{fig:reps:pretrained}
\end{subfigure}
\begin{subfigure}[t]{0.24\linewidth}
\centering
\begin{tikzpicture}[scale=0.6]
\begin{axis}[
 title={},
 grid=major,
axis lines=middle,
 no markers,
 every axis plot/.append style={ultra thick},
 xmin=-0.4999, xmax=0.55,
 ymin=-0.4999, ymax=0.55,
 ticks=none
]
\node[regular polygon, minimum height=6mm,regular polygon sides=3, fill=red!100,inner sep=2pt] at (axis cs:0.25644314, 0.12540874) {};
\node[regular polygon, minimum height=6mm,regular polygon sides=4, fill=blue!100] at (axis cs:0.25999054, -0.29644176) {};
\node[regular polygon, minimum height=6mm,regular polygon sides=5, fill=red!100] at (axis cs:0.049219638, 0.42767277) {};
\node[regular polygon, minimum height=6mm,regular polygon sides=6, fill=blue!100] at (axis cs:-0.061883904, 0.19825003) {};
\node[regular polygon, minimum height=\minshape,regular polygon sides=7, fill=red!100] at (axis cs:-0.37857392, 0.47314683) {};
\node[regular polygon, minimum height=\minshape,regular polygon sides=3, fill=blue!100,inner sep=2pt] at (axis cs:0.10887167, -0.26070255) {};
\node[regular polygon, minimum height=\minshape,regular polygon sides=4, fill=red!100] at (axis cs:-0.34162185, 0.050839007) {};
\node[regular polygon, minimum height=\minshape,regular polygon sides=5, fill=blue!100] at (axis cs:0.05225141, -0.4067908) {};
\node[regular polygon, minimum height=\minshape,regular polygon sides=6, fill=red!100] at (axis cs:0.46225715, 0.4129299) {};
\node[regular polygon, minimum height=\minshape,regular polygon sides=7, fill=blue!100] at (axis cs:-0.038552105, -0.38253385) {};
 %\node[regular polygon, minimum height=\minshape,regular polygon sides=3, fill=blue!100,inner sep=2pt] at (axis cs:-0.35932082, 0.17664503) {};
 %\node[regular polygon, minimum height=\minshape,regular polygon sides=3, fill=red!100,inner sep=2pt] at (axis cs: 0.23422007, -0.23213634) {};
 %\node[regular polygon, minimum height=\minshape,regular polygon sides=4, fill=blue!100] at (axis cs:0.4148342 , 0.02179685) {};
 %\node[regular polygon, minimum height=\minshape,regular polygon sides=4, fill=red!100] at (axis cs:-0.19870198, -0.02996263) {};
 %\node[regular polygon, minimum height=\minshape,regular polygon sides=5, fill=blue!100] at (axis cs: 0.23975432 , 0.40730593) {};
 %\node[regular polygon, minimum height=\minshape,regular polygon sides=5, fill=red!100] at (axis cs:0.49842966 , 0.17170991) {};
\end{axis}
\end{tikzpicture}
 \caption{The matrix in \figref{fig:reps:random}.}

 \label{fig:tsne-random}
\end{subfigure}
\begin{subfigure}[t]{0.24\linewidth}
\centering
\begin{tikzpicture}[scale=0.6]
\begin{axis}[
 title={},
 no markers,
 grid=major,
axis lines=middle,
 every axis plot/.append style={ultra thick},
 xmin=-0.9, xmax=1.8,
 ymin=-1.99, ymax=0.99,
ticks=none
]
\node[regular polygon, minimum height=6mm,regular polygon sides=3, fill=red!100,inner sep=2pt] at (axis cs:0.043389633, 0.8164206) {};
\node[regular polygon, minimum height=6mm,regular polygon sides=4, fill=blue!100] at (axis cs:-0.7053148, -1.812479) {};
\node[regular polygon, minimum height=6mm,regular polygon sides=5, fill=red!100] at (axis cs:-0.19551103, 0.4328718) {};
\node[regular polygon, minimum height=6mm,regular polygon sides=6, fill=blue!100] at (axis cs:1.6731963, -1.0493283) {};
\node[regular polygon, minimum height=\minshape,regular polygon sides=7, fill=red!100] at (axis cs:-0.19127868, -0.034780912) {};
\node[regular polygon, minimum height=\minshape,regular polygon sides=3, fill=blue!100,inner sep=2pt] at (axis cs:-0.40698338, -0.91700864) {};
\node[regular polygon, minimum height=\minshape,regular polygon sides=4, fill=red!100] at (axis cs:-0.24272422, 0.16861145) {};
\node[regular polygon, minimum height=\minshape,regular polygon sides=5, fill=blue!100] at (axis cs:0.88247913, -0.13587654) {};
\node[regular polygon, minimum height=\minshape,regular polygon sides=6, fill=red!100] at (axis cs:-0.075923756, 0.29488102) {};
\node[regular polygon, minimum height=\minshape,regular polygon sides=7, fill=blue!100] at (axis cs:2.1035564, -2.6820786) {};
 %\node[regular polygon,regular polygon sides=3, fill=blue!100,inner sep=2pt] at (axis cs:-1.2906312,  0.36732233]) {};
 %\node[regular polygon,regular polygon sides=3, fill=red!100,inner sep=2pt] at (axis cs: -0.09487347, -1.1454606) {};
 %\node[regular polygon,regular polygon sides=4, fill=blue!100] at (axis cs:-0.8569037 , 1.1844875) {};
 %\node[regular polygon,regular polygon sides=4, fill=red!100] at (axis cs: 0.56387985, -0.61090606) {};
 %\node[regular polygon,regular polygon sides=5, fill=blue!100] at (axis cs:-0.94791865, 2.0959575) {};
 %\node[regular polygon,regular polygon sides=5, fill=red!100] at (axis cs:1.2290089, -0.08361815) {};
\end{axis}
\end{tikzpicture}
 \caption{The matrix in \figref{fig:reps:pretrained}.}
 \label{fig:tsne-pretrained}
\end{subfigure}

\caption{Each matrix is a method for representing the shapes in \figref{fig:entities}, where each row is a vector representation of one shape. Localist and \newatticus{property} are featural representations where each vector dimension encodes the value of a single, \newmike{semantically interpretable} property. Random and pretrained are non-featural representations where the values of properties are encoded implicitly in two dimensions. Random representations and localist representations encode only identity, whereas \newatticus{property} representations and pretrained representations encode color and number of sides.}
\label{fig:reps}
\end{figure*}




\section{Designing theoretical models of equality learning}

We begin by discussing two critical design considerations for our models: (1) the standards for generalization by which models should be evaluated and (2) the type of representations they should use. To summarize this discussion: we select generalization tasks with fully disjoint training and test vocabularies to provide the most stringent test of generalization. Further, we adopt both randomly initialized  representations and pretrained representations for our subsequent models, and we show analytically that other kinds of representations are more limited in their capacity to make successful out-of-sample generalizations.

\subsection{Generalization}

The standard approach to training and evaluating neural networks is to choose a dataset, divide it randomly into training and assessment sets, train the system on the training set, and then use its performance on the assessment set as a proxy for its capacity to generalize to new data.

The standard approach is fine for many purposes, but it raises concerns in a context in which we are trying to determine whether a network has truly acquired a global solution to a target function. In particular, where there is any kind of overlap between the training and assessment vocabularies (primitive elements), we can't rule out that the network might be primarily taking advantage of idiosyncrasies in the underlying dataset to effectively cheat -- to memorize aspects of the training set and learn a local approximation of the target function that happens to provide traction during assessment.

To address this issue, \update{we follow \cite{marcus:1999} in proposing} that networks must be evaluated on assessment sets that are completely disjoint in every respect from the train set, all the way down to the entities involved. For example, below, we train on pairs $(a, a)$ and $(a, b)$, where $a$ and $b$ are representations from a train vocabulary $V_{T}$. At test time, we create a new assessment vocabulary $V_{A}$, derive equality and inequality pairs $(\alpha, \alpha)$ and $(\alpha, \beta)$ from that vocabulary, and assess the trained network on these new examples. In adopting these methods, we get a clear picture of the system's capacity to generalize, and we can safely say that its performance during assessment is a window into whether a global solution to identity has been learned. This is a very challenging setting for any machine learning model.
% For the sequential same--different task, we will see that it even requires us to depart from usual model formulations.

\subsection{Representations}\label{sec:representations}

% \updatea{
Essentially all modern machine learning models represent objects using vectors of real numbers. However, there are important differences in how these vectors are used to encode the properties of objects. \newatticus{The method of representation impacts whether there is a natural notion of similarity between entities and the ability of models to generalize to examples unseen in training. These two attributes are deeply related; if there is a natural notion of similarity between vector representations, then models can generalize to inputs with representations that are similar to those seen in training.}

We characterize two broad approaches to such property encoding -- which we call \tech{featural representations} and \tech{non-featural representations} -- and argue that the differences between them have not been given sufficient attention in the debate about the ability of neural networks to perform relational reasoning. We acknowledge that a dimension of any vector representation is a ``feature'' but we adopt a usage that is common in cognitive science, namely that a feature is an interpretable semantic primitive.\footnote{\newatticus{The term \tech{distributed representations} is used to refer jointly to what we call property representations, random representations, and pretrained representations. We opted not to use this term because it does not seperate property representations from random and pretrained representations, which is the relevant division here. Distributed representations are often contrasted in the neural network literature with \tech{local} or \tech{localist} representations; as discussed below, here we define these terms specifically to refer to representations whose features correspond to specific entities.}}

We ground our discussion in a hypothetical universe of blocks which vary by shape and color. \Figref{fig:entities} is a partial view of them, and \figref{fig:reps:localist}--\figref{fig:reps:pretrained} present four different ways of encoding the properties of these objects in vectors.

\subsubsection{Featural Representations}

The defining characteristic of \tech{featural} vector representations is that each dimension encodes the value of a single, \newmike{semantically intepretable} property. The properties can be binary, integer-valued, or real-valued.

We use the term \tech{localist} for the special case of featural representations in which only objects are represented and there is a feature corresponding to each object. In \figref{fig:reps:localist}, each column represents the property of being an object, and every object is represented as a vector that has a single unit with value 1. There is no shared structure across objects; all are equally (un)related to each other as far as the model is concerned.

\newatticus{We will refer to featural representation that are not localist as \tech{property representations}.} Here, column dimensions encode specific, meaningful properties of objects. In our example, we can represent the properties of being red and being blue with two different binary features, and the property of having a certain number of sides as a single integer feature, as in \figref{fig:reps:symbolic}. Unlike with localist representations, objects in this space can have complex relationships to each other, as encoded in the shared structure given by the columns.

Featural representations -- both localist and \newatticus{property} -- have the appealing property that they are easy for researchers to interpret because of the tight correspondence between column dimensions and properties. However, this transparency actually inhibits neural networks from discovering general solutions. Instead, such models work far better with representations that have property values implicitly encoded in the abstract structure of the vector space. We demonstrate this result analytically in \appref{app:generalization} for the case of binary features. The core insight is that networks cannot learn anything about column dimensions that are not represented in their training data; whatever weights are associated with those dimensions are unchanged by the learning process, so predictions about those dimensions remain random at test time. \newatticus{A developmental perspective suggests another reason to avoid binary representations, namely that this is not an accurate account of how perceptual inputs are represented in the brain.}

Recent work in machine learning \cite{weyde2019,weyde2018,kopparti2020} attempts to overcome this analytic limitation of binary featural representations by modifying standard neural architectures to have symbolic primitives or changing network weight priors. In our work, we instead opt for non-featural representations, which do not have this analytic limitation and are the norm in state-of-the-art artificial intelligence models. \newatticus{There is no need to introduce symbolic primitives or modify network weight priors when non-featural representations are used.}

\subsubsection{Non-Featural Representations}

A \tech{non-featural representation} is a vector that encodes property values implicitly across many dimensions. Perhaps the simplest non-featural representations are \tech{completely random} vectors, as in \figref{fig:reps:random}. Random representations can be seen as the non-featural counterpart to localist representations. In both of these representation schemes, all the objects are equally (un)related to each other, since column-wise patterns are unlikely in random representations and, to the extent that they are present, they exist completely by chance. However, in random representations, all the column dimensions can contribute meaningfully to identifying objects, whereas a localist representation has only one vector unit that determines the identity of any given object.

Random representations are a starting point that encodes object identity, but we can \tech{pretrain} these representations via a learning process, imbuing them with rich structure that implicitly encodes property values across many dimensions. \Figref{fig:reps:pretrained} provides a simple example. This matrix is the results of pretraining the representations in \figref{fig:reps:random} on the task of predicting whether the object is blue, whether the object is red, and the number of sides the object has. (\Appref{app:pretraining} provides technical details on our pretraining approach.) Superficially, the two matrices look equally random, but the random representations in \figref{fig:tsne-random} have no such structure, while the pretrained representations in \figref{fig:tsne-pretrained} do: there is a line that separates blue and red objects.

Pretraining need not be restricted to input representations; all the parameters of a model can be pretrained, offering the possibility that networks might be used as modular components to solve more complex tasks. We realize this possibility with our third experiment, where a model pretrained on a simple equality is used as a modular component to compute hierarchical equality.
% In this setting, a small amount of task-specific training (often called ``fine-tuning'') might suffice.

% \subsubsection{The Superiority of Non-Featural Representations}

% We argue that non-featural representations are superior to featural ones for two reasons: (1) non-featural representations are more biologically plausible than feature representations; and (3) featural representations can seriously inhibit the ability of neural networks to generalize to unseen examples.

% From a biological perspective, we might want the representations we feed our artificial neural networks to be similar to perceptual neural representations in biological neural networks. We believe it is obvious that biological neural representations are such that property values are not explicitly encoded or localized to the activity of a single neuron. For example, when a human looks at a red square block that is red, the property of being red is not encoded in the activity of a single neuron in the vision system, but instead is encoded in the activations of many neurons. This is a property that is shared with pretrained non-featural representations, but not with feature representations, where property values are explicit and localized.



% \subsubsection{On the Symbolic/Non-Symbolic Distinction}


% \subsection{The current paper}

%\begin{figure*}[tp]
%\centering
%\begin{subfigure}[t]{0.48\linewidth}
% \includegraphics[width=1\linewidth]{../fig/equality-train_size-embed_dim-hidden_dim=100.pdf}
% \vspace{-4mm}
% \caption{Results for a model in which the hidden layer has dimensionality $100$. The lines correspond to different dimensions for the entities $a$ and $b$ in the input pairs $(a, b)$. The best model passes 90\% accuracy with 1,000 training examples and reaches (near) perfection by 1,250.}
% \label{fig:equality--smallresults}
%\end{subfigure}
%\hfill
%\begin{subfigure}[t]{0.48\linewidth}
% \centering
% \includegraphics[width=1\linewidth]{../fig/equality-pretrain-compare-train_size-pretrained-embed_dim=None.pdf}
% \caption{\update{Results where random representations are grounded in a number of property domains via pretraining learning tasks. The `no pretraining' model is the best of the models at left (10-dimensional embeddings, 100-dimensional hidden layers), repeated to facilitate comparisons. The pretraining models use those same dimensionalities.}}
% \label{fig:basic-equality-pretrain}
%\end{subfigure}
%\caption{Same--different results with and without pretraining.}
%\end{figure*}

\section{Model 1: Same--different relations}\label{sec:equality}

First, we investigate whether a supervised single layer feed-forward neural network can learn the equality relation in the strict setting we describe above where train and test vocabulary are disjoint. The input is a pair of vectors $(a, b)$ which correspond to the two stimulus objects. These vectors are non-featural representations that do not have features encoding properties of the objects or their identity. During training, this model is presented with positive and negative labeled examples. During testing, this model is tasked with categorizing inputs unseen during training. It is straightforward to show that a network like this is capable of learning equality as we have defined it. Indeed, in our Supplementary Materials, we provide an analytic solution to the equality relation using this neural model.
% \Appref{app:equality-solution} provides

This result shows that equality in our sense is learnable in principle, but it doesn't resolve the question of whether networks can find this kind of solution given finite training data. To address this issue, we train networks on a stream of pairs of random vectors. Half of these are identity pairs $(a, a)$, labeled with $1$, and half are non-identity pairs $(a, b)$, labeled with $0$. Trained networks are assessed on the same kind of balanced dataset, with vectors that were never seen in training so that, as discussed earlier, we get a clear picture of whether they have found a generalizable solution.

\subsection{Results}

\Figref{fig:basic-equality-pretrain} shows our results. The representations used in these experiments are random representations that were pretrained using a linear classifier for 0, 3, 5, or 10 different binary feature discrimination tasks. For example, following \figref{fig:reps}, a three-task model might be trained to encode the binary properties of being blue, having four sides, and being red. For all representations, this neural model reached above-chance performance almost immediately, but required upwards of 1,000 examples to achieve near perfect accuracy.   Interestingly, we observed a clear speed-up, with more pretraining tasks resulting in the largest gains. It seems that, by grounding our representations in ``property domains'' (as represented by the different task dimensions), we imbue them with implicit structure that makes learning easier.

%\Figref{fig:equality--smallresults} presents typical results. This is for the case where the hidden layer dimensionality is $100$, and we plot results for different embedding dimensionalities $m$ and different amounts of training data. The picture is comparable with hidden dimensionalities at 10, 25, and 50, but those models require more training data to reach (near) perfect performance (\appref{app:model1-results}).

%While all the networks in \figref{fig:equality--smallresults} reach above-chance performance almost immediately, they require upwards of 1,000 examples to truly solve these tasks. We additionally ran the above experiments using random representations that were pretrained using a linear classifier for 0, 3, 5, or 10 different feature discrimination tasks. For example, following \figref{fig:reps}, a three-task model might be trained to encode the properties of color, number of sides, and size.

%\Figref{fig:basic-equality-pretrain} summarizes the results of these experiments for 10-dimensional embeddings and 100-dimensional hidden representations, as this seems to be the network that learns the fastest with random inputs. Interestingly, we see a clear speed-up, with more pretraining tasks resulting in the largest gains; by grounding our representations in ``property domains'' (as represented by the different task dimensions), we imbue them with implicit structure that makes learning easier.

\subsection{Discussion}

Our assessment pairs have nothing in common with the training pairs except insofar as both involve vectors of real numbers of the same dimensionality. During training, the network is told (via labels) which pairs are equality pairs and which are not, but the pairs themselves contain no information about equality per se. It thus seems fair to us to say that these networks have learned equality -- or at least how to simulate that relation with near perfect accuracy. Further, the use of representations that are structured by pretraining results in faster learning.



%\begin{figure*}[tp]
%\centering
%\begin{subfigure}[t]{0.48\textwidth}
% \centering
% \includegraphics[width=1\linewidth]{../fig/fuzzy-lm-vocab20-train_size-embed_dim-hidden_dim=100.pdf}
% \caption{Results for a model in which each $h_{t}$ in Eq.~\eg{eq:lstm-recur} has dimensionality 100. The lines correspond to the dimensionality of the input representations ($x_t$ in Eq.~\eg{eq:lstm-recur}). All the training examples are presented at once over multiple epochs.}
% \label{fig:fuzzy-lm-results}
%\end{subfigure}
%\hfill
%\begin{subfigure}[t]{0.48\textwidth}
% \centering
% \includegraphics[width=1\linewidth]{../fig/fuzzy-lm-pretrain-compare-train_size-pretrained-embed_dim=None.pdf}
% \caption{Results where random representations are grounded via pretraining learning tasks. Simulations adopt the best-performing settings from the no-pretraining condition (2-dimensional embeddings, 100-dimensional hidden representations).}
% \label{fig:fuzzy-lm-pretrain-results}
%\end{subfigure}
%\caption{Sequential same--different results with and without pretraining.}
%\end{figure*}


\begin{figure*}[tp]
\centering
\begin{subfigure}[t]{0.48\linewidth}
 \centering
 \includegraphics[width=1\linewidth]{../fig/equality-pretrain-compare-train_size-pretrained-embed_dim=None.pdf}
 \caption{Results for single layer feed-forward neural networks trained on our simple equality task. The 'no pretraining' model is provided random representations and the 'k-task pretraining' models are provided random representations are grounded in $k$ binary property domains via pretraining learning tasks. }
 \label{fig:basic-equality-pretrain}
\end{subfigure}
\hfill
\begin{subfigure}[t]{0.48\textwidth}
 \centering
 \includegraphics[width=1\linewidth]{../fig/fuzzy-lm-pretrain-compare-train_size-pretrained-embed_dim=None.pdf}
 \caption{Results for LSTM recursive neural networks trained on our sequential equality task. The 'no pretraining' model is provided random representations and the 'k-task pretraining' models are provided random representations are grounded in $k$ binary property domains via pretraining learning tasks.  All the training examples are presented at once over multiple epochs}
 \label{fig:fuzzy-lm-pretrain-results}
\end{subfigure}
\end{figure*}





\section{Model 2: Sequential same--different (ABA task)}\label{sec:lms}

Our first model is simple and successfully learns equality. However, this model is supervised with both positive and negative evidence. In the initial debate around these issues, supervision with negative evidence was dismissed as an unreasonably strong learning regime \citep[e.g.,][]{marcus:1999a}. While this argument likely holds true for language learning \citep[in which supervision is generally agreed not to be binary or direct;][]{brown:1970,chouinard2003}, it is not necessarily true for learning more generally. Nevertheless, learning of sequential rules without negative feedback is possible for infants \cite{marcus:1999,rabagliati:2019}. In experiments of this type, infants are presented with a set of positive examples. Our next model explores whether neural network models can learn this task in a challenging regime with no negative supervision.

To explore learning with only positive instances, we use a neural LSTM language model, a recursive network with the ability to selectively forget and remember information \cite{hochreiter:1997}. Language models are sequential: at each timestep, they predict an output given their predictions about the preceding timesteps. As typically formulated, the prediction function is just a classifier: at each timestep, it predicts a probability distribution over the entire vocabulary of options, and the item with the highest probability is chosen as a symbolic output. This output becomes the input at the next timestep, and the process continues.

This formulation will not work in situations in which we want to make predictions about test items with an entirely disjoint vocabulary from the training sample. The classifier function will get no feedback about these out-of-vocabulary items during training, and so it will never predict them during testing. To address this issue, we reformulate the prediction function. Our proposal is to have the model predict output vector representations -- instead of discrete vocabulary items -- at each timestep. During training, the model is trained to minimize the distance between these output predictions and the representations of the actual output entities. During assessment, we take the prediction to be the item in the entire vocabulary (training and assessment) whose representation is closest to the predicted vector (in terms of Euclidean distance). This fuzzy approach to prediction creates enough space for the model to predict sequences from an entirely new vocabulary. Our Supplementary Materials
% \Appref{sec:analyticlm}
provide an analytic solution to the ABA task using this model.

To see how well the model performs in practice, we trained networks on sequences \texttt{<s> a b a </s>}, where $\texttt{b} \neq \texttt{a}$. We show the network every such sequence during training, from an underlying vocabulary of 20 items (creating a total of 380 examples). To assess how well the model learns this pattern, we seed it with \texttt{<s> x} where \texttt{x} is an item from a disjoint vocabulary from that seen in training, and we say that a prediction is accurate if the model continues with \texttt{y x </s>}, where \texttt{y} is any character (from the training or assessment vocabulary) except \texttt{x}.

\subsection{Results}

\Figref{fig:fuzzy-lm-pretrain-results} shows our results. Unlike for the previous equality experiment, we found that we had to allow the model to experience multiple epochs of training on the same set in order to succeed and tens of thousands of training examples were necessary. We considered a range of representations (as in Model~1); the model was again successful with all representations, but in this experiment pretraining representations did not increase performance.


\subsection{Discussion}

These sequential models are given no negative examples and they must predict into a totally new vocabulary. Despite these challenges, they succeed at learning the underlying patterns in our data. On the other hand, the learning process is slow and data-intensive. We hypothesized that grounding representations in property domains via pretraining might lead to noticeable speed-ups, as it did in for our simple same-different task, but we did not see this effect in practice. We speculate that there may be model variants that reduce these demands, given that learning is in principle possible in this architecture, but we leave them to future work.

\section{Model 3: Hierarchical same--different relations}\label{sec:premack}

Given the strong results found for simple equality relations, we can ask whether more challenging equality problems are also learnable in our setting. The hierarchical equality task used by Premack \cite{Premack:1983} is an interesting test case: given a pairs of pairs $((a,b), (c,d))$, the label is $1$ if $(a = b) = (c = d)$, else $0$. Premack suggested that the ability exemplified by this task -- reasoning about hierarchical \emph{same} and \emph{different} relations -- could represent a form of symbolic abstraction uniquely enabled by language. Given the non-symbolic nature of our models, our simulations provide a test of this hypothesis, though we should look critically at their ability to find good solutions with reasonable amounts of training data.


\begin{figure*}[tp]
\centering
\begin{subfigure}[t]{0.48\linewidth}
 \centering
 \includegraphics[width=1\linewidth]{../fig/flatpremack-h2-pretrain-compare-train_size-pretrained-embed_dim=None.pdf}
 \caption{Results for two layer feed-forward neural networks trained on our hierarchical equality task. The 'no pretraining' model is provided random representations and the 'k-task pretraining' models are provided random representations are grounded in $k$ binary property domains via pretraining learning tasks. }
 \label{fig:premack-h2-pretrain}
\end{subfigure}
\hfill
\begin{subfigure}[t]{0.48\textwidth}
\includegraphics[width=1\linewidth]{../fig/input-as-output-pretrain-compare-train_size-pretrained-embed_dim=None.pdf}
\caption{Results for simple equality networks applied to the hierarchical equality task. The 'no pretraining' model is provided random representations and the 'k-task pretraining' models are provided random representations are grounded in $k$ binary property domains via pretraining learning tasks.
%Even with no additional training instances for this task, all models achieve greater than chance accuracy, and even modest amounts of additional training on the task lead to excellent performance.
}
\label{fig:premack-pretraining-results}
\end{subfigure}
\end{figure*}
%\begin{figure*}[tp]
%\centering
%\begin{subfigure}[t]{0.48\linewidth}
% \centering
% \includegraphics[width=1\linewidth]{../fig/flatpremack-h2-train_size-embed_dim-hidden_dim=100.pdf}
% \caption{Results for a network with two 100-dimensional hidden layers. Nearly all the networks solve the task, but they require very large training sets to do so.}
% \label{fig:premack-h2-flat-results}
%\end{subfigure}
%\hfill
%\begin{subfigure}[t]{0.48\linewidth}
% \centering
% \includegraphics[width=1\linewidth]{../fig/flatpremack-h2-pretrain-compare-train_size-pretrained-embed_dim=None.pdf}
% \caption{Results from pre-trained networks, adopting the best of the models at left (10-dimensional embeddings, 100-dimensional hidden representations).}
% \label{fig:premack-h2-pretrain}
%\end{subfigure}
%\caption{Hierarchical same--different results with and without pretraining.}
%\end{figure*}

We can approach this task using the same model and methods as we used for equality, with the relatively minor change of providing the network four vector representations instead of two. We found that single layered feed-forward neural networks required nearly 100,000 training examples to solve this task.  We hypothesized that the flat input representations $[a;b;c;d]$ might be suboptimal here. This task is intuitively hierarchical: if one works out the equality labels for each of the two pairs, then the further classification decision can be done entirely on that basis. Our current neural network might be too shallow to find this kind of decomposition. To address this issue, we use a two layer feed forward network.

\subsection{Results without pretraining}

\Figref{fig:premack-h2-pretrain} shows our results. We again considered a range of representations and again the network succeeded across this range, with pretraining increasing performance as in Model~1. The network required more than 20,000 training instances to reach top performance, however, and upwards of 10,000 examples with pretrained representations.

This amount of training data is vastly more data than human participants get in similar experiments, which typically involve short exposures in the range of dozens to hundreds of examples \cite[e.g.,][]{marcus:1999,endress2005}. Thus, it is worth asking whether there are other solutions that would be more data efficient and more in line with human capabilities. We next seek to further capitalize on the hierarchical nature of this task by defining a modular pretraining regime in which previously learned capabilities are recruited for new tasks.

\subsection{The critical role of experience}\label{sec:modular}

Our successful results training neural networks on simple equality suggested another strategy for solving the hierarchical equality task. Rather than requiring our networks to find solutions from scratch, we pretrained them on basic equality tasks and then used those parameters as a starting point for learning hierarchical equality. This set of simulations was conceptually similar to our previous experiments with pretrained input representations, but now we pretrained an entire subpart of the model, rather than just input representations.

The hierarchical equality task requires computing the equality relation three times: compute whether each pair of inputs are equal and then compute whether the truth-valued outputs of these first two computations are equal. We thus used the same network pretrained on basic equality to perform all three equality computations.

%\begin{figure*}[tp]
%\centering
%\begin{subfigure}[t]{0.48\textwidth}
%\includegraphics[width=1\linewidth]{../fig/input-as-output-train_size-embed_dim-hidden_dim=0.pdf}
%\caption{Basic equality networks applied to the hierarchical equality task. Each line represents a different embedding dimensionality, which is constrained in this model to match the hidden dimensionality. Even with no additional training instances for this task, all models achieve greater than chance accuracy, and even modest amounts of additional training on the task lead to excellent performance.
%}
%\label{fig:premack-pretraining-results}
%\end{subfigure}
%\hfill
%\begin{subfigure}[t]{0.48\textwidth}
% \centering
% \includegraphics[width=1\linewidth]{../fig/input-as-output-zero-shot.pdf}
% \caption{\update{The results from the plot at left where the x-axis is 0, that is, where we are testing zero-shot generalization from the pretrained equality network to the hierarchical setting. All the models perform well above chance in this setting, with the 50-dimensional version achieving a mean of 67\% accuracy.}}
% \label{fig:premack-pretraining-zero-shot-results}
%\end{subfigure}
%\caption{Modular network results for the hierarchical same--different task.}
%\label{fig:modular-results}
%\end{figure*}

% \subsubsection{Results and discussion}

\Figref{fig:premack-pretraining-results} shows our results. All the models have above chance performance after being trained only on the simple equality task -- that is, they achieve zero-shot generalization to the hierarchical task and within two thousand examples, the models achieve near perfect accuracy. It is remarkable that a model trained only on equality between entities is able to get traction on a problem that requires determining whether equality holds between the truth values encoded in two learned representations. Pretrained representations did not increase performance.

%It might be possible to effectively combine network pretraining with input pretraining as in the previous experiments. An initial exploration of this idea is presented in \appref{app:double-pretrain}. While we have not yet found a way to use this combination of pretraining regimes to improve over \figref{fig:modular-results}, we are optimistic about such combinations for future work.

\section{General Discussion}

% Further, an important line of work on relational reasoning in neural network architectures sidesteps the issue of whether relational
 % the issue of whether equality relations can be induced from data by creating networks whose architecture is explicitly relational \citep{raposo2017,santoro2018,santoro2019,palm2019}.

Equality is a key case study for understanding the origins of human relational reasoning. This case study has been puzzling for symbolic accounts of reasoning because such accounts do not provide a compelling explanation for why some equality tasks are so easy to learn and others are so hard. In addition, evidence of graded learning and generalization in non-human species suggests that a gradual learning account might provide more traction in explaining the empirical data \cite{wasserman:2017}. Inspired by this work, we revisited a long-standing debate about whether neural network models can learn equality relations from data \cite{alhama:2019}.

\newmike{We make three contributions to this debate. First, we show that non-featural representations –- both random and  pretrained\footnote{\newatticus{Pretrained representations were not a necessary component for models to succeed, and only led to learning speed increases for the simple and hierarchical equality tasks.}} -- allow standard neural networks to learn simple, sequential, and hierchical equality tasks. Both the research that originated this debate \cite{marcus:1999,marcus:2001,dienes:1999,seidenberg:1999a,seidenberg:1999b,elman:1999,negishi:1999} and more recent work \cite{weyde2019, weyde2018, kopparti2020, alhama2018} only involve experiments where featural representations are used, and we suggest that this choice led directly to the conclusions of this body of work. Second, we show that neural networks can achieve high test accuracy on the sequential equality task with no negative feedback, suggesting that a negative feedback learning regime is not critical for learning equality. Finally, we show that a neural network trained only on simple equality can generalize to hierarchical equality, even in a ``zero-shot'' evaluation.}

% We further supported these results with analytic insights into why networks are able to find solutions to these tasks.

In some settings, our current models require many more training instances than humans seem to need. However, our pretraining approach suggests a path forward: by using pretrained models as modular components, we can get traction on challenging tasks without any training specifically for those tasks. In some cases, even a small amount of additional training can make a substantial difference. Perhaps pretrained components of this type could serve as the basis for more complex cognitive abilities more generally \cite{frank2008,heyes2018}. \newmike{Other computational work on relational reasoning in cognitive science has explored hybrid architectures that do not explicitly encode symbolic equality but incorporate other symbolic structures \cite[e.g.,][]{hummel:2003}; integrating functions learned from data could be an interesting direction for future work with these architectures. Similarly, other work on relational reasoning in artificial intelligence has explored networks with explicitly relational architectures -- such architectures could be interesting to explore in the context of equality reasoning \citep{raposo2017,santoro2018,santoro2019,palm2019}.}

One further implication of our pretraining findings is that it should be possible to scaffold non-human animals' performance in complex, hierarchical equality tasks via training on simpler ones. Indeed, \cite{smirnova2015} show just this result in crows, consistent with our findings. Although we do not discount the potential role of linguistic labels in informing adult humans' expertise in such tasks \cite{gentner2003}, pretraining also provides a potential account of how infants and young children might succeed in a range of equality reasoning tasks without access to specific linguistic symbols like ``same'' \cite{walker:2016,ferry:2015,hochmann2016}.

More broadly still, our work suggests a possible way forward in understanding the acquisition of logical semantics. Graded logical functions like those our models learned here could form the foundation for a semantics of words like ``same'' \cite{potts2019}. Such an option is appealing because it escapes from the circularity of defining the semantics of linguistic symbols as originating in a mental primitive {\sc same}. A semantics for ``same'' requires defining its inputs and outputs as well as how it composes with other symbols. The assertion that there is a primitive identity computation does not specify the format of these inputs and outputs or these composition rules; it further fails to explain the flexibility that allows us to call two Toyota Corollas ``the same'' but two twin sisters ``different.'' In contrast, the kinds of networks we propose here could in principle be conditioned contextually to provide flexible, context-sensitive interpretation of logical meaning. \newmike{Such a contextually-conditioned semantics could in principle be applied in both cases of sameness and graded similarity \cite{medin1993}, holding the promise of unifying models of identity and similarity.}   % MCF: could we cite some of Ruben's work or the natural language inference work here maybe? or is there something more appropriate?

Earlier debates about the nature of equality computations centered around the question of whether models included symbolic elements. We believe ours do not; but it is of course possible to quibble with this judgment. For example, since the supervisory signal used in Models 1 and 3 is generated based on a symbolic rule, perhaps that makes these models symbolic under some definition. (Of course, the same argument could be applied to the supervision signal that is provided to crows, baboons, and human children in some tasks). We view this kind of argument as terminological, rather than substantive. In the end, our goal is an explicit learning theory for relational reasoning. Our hope is that the work described here takes a first step in this direction.

% The claim that certain forms of representation are equivalent to including symbolic primitives is untroubling if these representations enhance the generalization skills of neural models; some class of neural models might be labeled as symbolic, but that does not nullify the fact that they are successful models of graded relational learning and distinct from uncontroversially symbolic models \cite[e.g.,][]{frank:2011}.


% BACK TO SYMBOLIC, LANGUAGE, SEMANTICS

% the debate has been whether there is a symbol. We view this as the least important issue. A single ``same" symbol cannot account for the vast range of behaviors that we want to explain -- how are two toyota corollas the same, while two people are not?
% To be an explanatory account, a symbol needs a semantics - what are the inputs, outputs, how does it compose. Without these, there is no account of behavior.

% same isn't a relation, it's a set of possible relations


% you can think of it this way: our networks are a semantics for the symbol

% This is maybe unfamiliar, because you might think you'd do this in a logical language or have an axiom system, but our networks are perfect and so you'd never notice. MARR cite



% We report results with and without pretraining. While \citeauthor{marcus:2001} might consider our pretraining method to be symbolic, we believe it would be difficult to argue that random distributed representations are symbolic, as they do not encode any properties directly. We additionally believe that the symbolic/non-symbolic distinction does not provide an explanation for the state of the comparative and developmental literature on relational learning, as we will argue in \secref{sec:?}.}

% SYMBOLIC VS. NON-Symbolic
% - this distinction doesn't helps
% - neural models clearly different
% - supervision unclear
% - not helpful (localist could be symbolic & unstructured)
%
% Since the initial debate, neural networks have become

%Finally, we note that, while the pretrained representations in \figref{fig:reps:pretrained} encode property values via linear structures, this is only because we used linear classifiers in our pretraining tasks. There are a practically limitless number of different ways to encode property values via abstract structures.

 % By contrast, the literature on relational learning in neural networks is currently limited to feature representations, and does not consider any non-featural representations where property values are encoded implicitly with abstract structures. For example, all of the experiments reviewed in \cite{alhama:2019} use feature representations, meaning that each vector unit explicitly contains the value of a property. Their negative conclusions about relational learning in neural networks hold only for networks using such representations.


\ifthenelse{
  \equal{\papermode}{pnas}
}{
  \matmethods{
The code and data used to run the experiments in this paper are publicly available at: https://github.com/atticusg/NeuralRelationalReasoning

The simple equality model takes in input vectors, $a,b$, and uses a single linear transformation followed by a non-linearity to create a hidden representation $h$
that is then used to create a probability distribution, $y$, over the two classes.
\begin{align}
  h &= \ReLU([a;b]W_{xh} + b_{h})  & y = \softmax(hW_{hy} + b_{y}) \label{eq:h2y}
\end{align}

The sequential equality model takes in a sequence of input vectors, $x_1,x_2,\dots$, and uses an LSTM cell to create a hidden representation $h_t$ at each timestep $t$ that is
that is then linearly projected into the input vector space providing a prediction for that timestep, $y_t$.
\begin{align}
  h_{t} &= \LSTM(x_{t}, h_{t-1})  & y_{t} = h_{t}W + b\label{eq:lstm-predict}
\end{align}

The first hierarchical equality model takes in input vectors, $a,b,c,d$, and applies a linear transformation followed by a non-linearity to create a hidden representation $h_1$ and then
applies these two steps once more to create second hidden representation $h_2$ that is then used to create a probability distribution, $y$, over the two classes.
\begin{align}
  h_{1} &= \ReLU([a;b;c;d]W_{xh} + b_{h_{1}}) \label{eq:x2h1}\\
  h_{2} &= \ReLU(h_{1}W_{hh} + b_{h_{2}}) & \hspace{-10pt}y = \softmax(h_{2}W_{hy} + b_{y}) \label{eq:h2y2}
\end{align}

The second hierarchical equality model takes in input vectors, $a,b,c,d$, and applies the simple equality model from equations \eg{eq:h2y} to the pairs $(a,b)$
and $(c,d)$ to produce hidden representations $h_1$ and $h_2$. Then the simple equality model is applied once again to the pair $(h_1,h_2)$ to produce a final hidden representation $h_3$
that is used to create a probability distribution, $y$, over the two classes.
\begin{align}
  h_1 &= \ReLU([a;b]W_{xh} + b_{h}) & h_2 = \ReLU([c;d]W_{xh} + b_{h})\\
  h_3 &= \ReLU([h_1;h_2]W_{xh} + b_{h}) & y = \softmax(h_3W_{hy} + b_{y}) \label{eq:prey}\
\end{align}

The parameters for the simple and hierarchical equality models are learned using back propagation with a cross entropy objective function defined as follows, for a set of $N$ examples and $K$ classes:
%
\begin{equation}\label{eq:crossent}
\max(\theta)
\quad
\frac{1}{N}
\sum_{i=1}^{N}
\sum_{k=1}^{K}
y^{i,k} \log(h_{\theta}(i)^{k})
\end{equation}
%
where $\theta$ abbreviates the model parameters, $y^{i,k}$ is the actual label for example $i$ and class $k$, and $h_{\theta}(i)^{k}$ is the corresponding prediction.

The parameters for the sequential model are learned using back propagation with a squared mean error objective function defined as follows:
%
\begin{equation}
\max(\theta)
\quad
-\frac{1}{N}
\sum_{i=1}^{N}
\sum_{t=1}^{T_{i}}
\left\| h_{\theta}\left(x^{i, 0:{t-1}}\right) - x^{i,t} \right\|^{2}
\end{equation}
%
for $N$ examples. Here, $T_{i}$ is the length of example $i$. As before, $\theta$ abbreviates the parameters of the model. We use $h_{\theta}(x^{i, 0:{t-1}})$ for the vector predicted by the model for example~$i$ at timestep~$t$, which is compared to the actual vector at timestep $t$ via squared Euclidean distance.

Hyperparameter searches and implementation details can be found in \appref{app:optimization}.

}

\showmatmethods{} % Display the Materials and Methods section

\acknow{\input{acknowledgements}}

\showacknow{} % Display the acknowledgments section

% Bibliography
\bibliography{relational-learning-bib}

\end{document}

}{
  \input{footer-arxiv}
}
% \update{

% As noted above, this is an unusual formulation for a language model. The usual version essentially treats every timestep as involving a classification decision, with a cross-entropy loss. We cannot adopt this because of our goal of using unseen vocabulary items at test time. % MCF: seemed redundant
